
\documentclass{article}
\usepackage{colortbl}
\usepackage{makecell}
\usepackage{multirow}
\usepackage{supertabular}

\begin{document}

\newcounter{utterance}

\centering \large Interaction Transcript for game `dond', experiment `coop\_en', episode 1 with selfrefine{-}Qwen3.5{-}9B.
\vspace{24pt}

{ \footnotesize  \setcounter{utterance}{1}
\setlength{\tabcolsep}{0pt}
\begin{supertabular}{c@{$\;$}|p{.15\linewidth}@{}p{.15\linewidth}p{.15\linewidth}p{.15\linewidth}p{.15\linewidth}p{.15\linewidth}}
    \# & $\;$A & \multicolumn{4}{c}{Game Master} & $\;\:$B\\
    \hline

    \theutterance \stepcounter{utterance}  
    & & \multicolumn{4}{p{0.6\linewidth}}{
        \cellcolor[rgb]{0.9,0.9,0.9}{
            \makecell[{{p{\linewidth}}}]{
                \texttt{\tiny{[P1$\langle$GM]}}
                \texttt{You are playing a cooperative negotiation game in which you have to agree on how to divide a set of items among you and another player.} \\
\\ 
\texttt{Rules:} \\
\texttt{(a) You and the other player are given a set of items. Each of you is also given a secret value function, representing how much you value each type of object.} \\
\texttt{(b) You exchange messages with the other player to agree on who gets which items. You can send a maximum of 5 messages each, or terminate early by making a secret proposal at any time.} \\
\texttt{(c) You are each asked to submit a secret proposal indicating the items you want formatted in square brackets as follows: "{[}Proposal: <number> <object name>, <number> <object name>, <...>{]}"} \\
\texttt{(d) If your proposals are complementary, i.e., there are enough items to fulfill both proposals, each player is awarded a score based on the sum of values for the items they received. Otherwise, both of you get zero points.} \\
\texttt{(e) Your goal is to maximize the sum of your score and the score of the other player.} \\
\\ 
\texttt{Let us start.} \\
\\ 
\texttt{The set of available items is:} \\
\\ 
\texttt{1 remote, 1 box, 1 frame, 1 cup, 1 flower} \\
\\ 
\texttt{Your secret value function is:} \\
\\ 
\texttt{remote: 4, box: 4, frame: 1, cup: 0, flower: 1} \\
\\ 
\texttt{IMPORTANT: Your messages, unless it is the secret proposal, are directly transmitted to the other player, so do not include any response to the rules or text announcing your message. To make a secret proposal, use the indicated format. Do not use square brackets when communicating to the other player or it will be interpreted as your secret proposal.} \\
            }
        }
    }
    & \\ \\

    \theutterance \stepcounter{utterance}  
    & \multicolumn{4}{p{0.6\linewidth}}{
        \cellcolor[rgb]{0.95,0.95,0.95}{
            \makecell[{{p{\linewidth}}}]{
                \texttt{\tiny{[P2$\rightarrow$P1]}}
                \texttt{{[}DRAFT{]}} \\
\texttt{I value the remote (4) and box (4) most, totaling 8. The frame and flower are worth 1 each, and the cup is worth 0. To maximize our combined score, we should aim for a division where we both get high{-}value items. Since there are 5 items and we need to split them, a 2{-}2{-}1 split isn't possible for two people; it must be a 3{-}2 split or similar.} \\
\\ 
\texttt{If I take the remote and box (value 8), I leave the frame, cup, and flower for you. You would value the frame (1) and flower (1) for a total of 2 (cup is 0). Total combined score: 10.} \\
\\ 
\texttt{Alternatively, if we split the high{-}value items differently, say I take the remote (4) and you take the box (4), we still have the frame, cup, and flower left. We could split those. But since I value the box equally to the remote, and you likely value them similarly (assuming symmetric or standard values for this type of puzzle unless specified otherwise, though I only know my own), the most efficient way to ensure we both get something valuable without conflict is for me to claim the two highest{-}value items I know of, leaving the rest for you to pick from, or proposing a specific split.} \\
\\ 
\texttt{However, the rule says "If your proposals are complementary...". This implies I should propose a specific set of items I want, and you propose a set you want. If they don't overlap and cover all items (or at least the items requested are available), we succeed.} \\
\\ 
\texttt{My strategy: I will propose to take the **remote** and the **box**. These are my highest value items. I will leave the frame, cup, and flower for you. This assumes you will propose to take those remaining items. If you propose to take the remote or box as well, we fail. But since I am maximizing the *sum* of our scores, and I know my values, I should propose the items that give me the most utility while leaving enough for you to make a non{-}zero proposal.} \\
\\ 
\texttt{Actually, looking at the goal: "maximize the sum of your score and the score of the other player."} \\
\texttt{My values: Remote(4), Box(4), Frame(1), Cup(0), Flower(1).} \\
\texttt{Total value to me: 10.} \\
\texttt{If I take Remote and Box, I get 8. You get Frame, Cup, Flower. If your values are similar (e.g., Box=4, Remote=4, Frame=1, Flower=1), you get 2. Sum = 10.} \\
\texttt{If I take Remote, Box, and Frame, I get 9. You get Cup, Flower. You get 1. Sum = 10.} \\
\texttt{If I take Remote, Box, Frame} \\
            }
        }
    }
    & & \\ \\

    \theutterance \stepcounter{utterance}  
    & \multicolumn{4}{p{0.6\linewidth}}{
        \cellcolor[rgb]{0.95,0.95,0.95}{
            \makecell[{{p{\linewidth}}}]{
                \texttt{\tiny{[P2$\rightarrow$P1]}}
                \texttt{{[}SELF{-}FEEDBACK{]}} \\
\texttt{FORMAT: problem} \\
\texttt{RULES: problem} \\
\texttt{STRATEGY: problem} \\
\texttt{REASON: The response is a long chain of thought containing internal monologue and incomplete sentences rather than a valid game message or secret proposal in the required format.} \\
            }
        }
    }
    & & \\ \\

    \theutterance \stepcounter{utterance}  
    & \multicolumn{4}{p{0.6\linewidth}}{
        \cellcolor[rgb]{0.9,0.9,0.9}{
            \makecell[{{p{\linewidth}}}]{
                \texttt{\tiny{[P1$\rangle$GM]}}
                \texttt{problem} \\
\texttt{The current thought process is too verbose and contains internal monologue rather than a direct game message. It fails to adhere to the constraint that non{-}secret messages must be transmitted directly to the other player without extra text, and it does not finalize a decision into the required secret proposal format.} \\
\\ 
\texttt{RULES} \\
\texttt{1. Messages must be direct transmissions to the other player without conversational filler.} \\
\texttt{2. Secret proposals must strictly follow the format: "{[}Proposal: <number> <object name>, <number> <object name>, <...>{]}"} \\
\texttt{3. The goal is to maximize the sum of both players' scores by ensuring complementary proposals.} \\
\texttt{4. The user has 5 messages available; this is the first move.} \\
\\ 
\texttt{STRATEGY} \\
\texttt{I need to propose a specific set of items that maximizes my value while leaving a valid, non{-}conflicting set for the other player. My highest value items are the remote (4) and the box (4). If I propose these two, I secure 8 points. The remaining items (frame, cup, flower) are worth 1, 0, and 1 respectively. Assuming the other player values items similarly or reasonably, they can take the frame and flower for a positive score, ensuring the combined score is maximized without overlap. Proposing all three high{-}value items might leave the other player with only the cup (value 0), potentially reducing the total sum if their valuation of the cup is also low or if they need at least one item to make a non{-}zero proposal. A 2{-}item proposal for me and a 2{-}item proposal for them seems optimal.} \\
\\ 
\texttt{REASON} \\
\texttt{The previous output was a stream of consciousness analyzing the game state, which violates the rule that messages must be direct transmissions. It also failed to commit to a specific proposal format. To proceed correctly, I must select the items I want (remote and box) and format them exactly as required for a secret proposal, assuming the other player will logically take the remaining items to avoid conflict and maximize the joint score. This move secures my highest value items while ensuring the game can succeed.} \\
\\ 
\texttt{{[}Proposal: 1 remote, 1 box{]}} \\
            }
        }
    }
    & & \\ \\

    \theutterance \stepcounter{utterance}  
    & & & \multicolumn{2}{p{0.3\linewidth}}{
        \cellcolor[rgb]{0.9,0.9,0.9}{
            \makecell[{{p{\linewidth}}}]{
                \texttt{\tiny{[GM$|$GM]}}
                \texttt{proposal does not include number/name pairs} \\
            }
        }
    }
    & & \\ \\

\end{supertabular}
}

\end{document}
